\pdfbookmark[1]{Abstract}{Abstract}
\chapter*{Abstract}

Biological systems exhibit complex, multi-scale dynamics spanning molecular interactions (milliseconds) to cellular processes (hours). Integrating these diverse timescales and regulatory mechanisms within a unified computational framework remains a fundamental challenge in systems biology.

This thesis presents the \textbf{Extended Biological Petri Net (Bio-PN)} formalism, introducing four core innovations:

\begin{enumerate}
    \item \textbf{Weak Independence Theory}: A relaxed form of transition independence permitting shared catalysts and convergent pathways while preserving reachability properties. Empirically, 64\% of biological transition pairs exhibit weak independence, enabling up to 3× parallel execution speedup.
    
    \item \textbf{Heterogeneous Transition Types}: Unified formal semantics for four dynamics---continuous (ODE), stochastic (Gillespie SSA), timed (scheduled events), and burst (transcriptional bursting)---coordinated through a hybrid scheduler.
    
    \item \textbf{Arc-Level Regulation}: Test arcs (catalysis) and inhibitor arcs (feedback) with threshold functions, providing compositional, topologically visible regulatory logic superior to global event systems.
    
    \item \textbf{Atomic Conservation}: Automatic verification of elemental balance via biochemical formula tracking, with database-driven cofactor suggestion (KEGG, ChEBI integration).
\end{enumerate}

The formalism is validated through 16 progressive examples and three comprehensive case studies: complete glycolysis (10 transitions, 3 regulatory checkpoints), citric acid cycle (8 transitions, cyclic topology), and integrated cellular respiration (32 transitions, carbon and energy accounting). All models exhibit physiologically realistic steady-state concentrations and regulatory responses.

Implementation in the SHYpn platform demonstrates 85--95\% modeling time reduction through automatic parameter inference from BRENDA (kinetic database). Performance benchmarks show linear time scaling (0.58s per transition) and 3× parallel speedup on 8 cores for weakly independent transitions. The system is competitive with COPASI (1.3× faster) and Snoopy (2× faster) despite Python overhead.

This work advances both Petri net theory (generalizing independence for catalytic systems) and systems biology practice (structured, automated, verifiable modeling). Future directions include spatial dynamics (colored Petri nets), genome-scale models (hierarchical abstraction, GPU acceleration), and machine learning integration (structure and parameter inference from omics data).

\vspace{1cm}
\noindent\textbf{Keywords:} Biological Petri nets, Multi-scale modeling, Weak independence, Hybrid simulation, Systems biology, Parameter inference
